\lecture{4. The Covenant I Made With Their Fathers}{04}

\section{Introduction}

\begin{frame}
\frametitle{Commitments can be scary}
\begin{center}
\includegraphics[trim = 1.5in .5in 1.5in .5in, clip=true, height=0.8\textheight]{figures/ballAndChain.jpg}
\end{center}

\note{09:30}
\note[item]{Being married can certainly cramp your style.}
\note[item]{You're accountable to someone for where you go, what you do, how you spend your money, who your friends are, etc.}
\note[item]{But, commitments also bring peace because it gives you someone you can count on}
\note[item]{Commitments make you a team.}
\note[item]{Commitments are the foundation of good relationships.}
\note[item]{Today we're going to dive into the commitments between us and God}
\note[item]{We're going to start talking about covenants}

\end{frame}

\begin{frame}
\frametitle{God wants a relationship with people}
\framesubtitle{Jeremiah 31:31-34}
	\keyversehiglight{the covenant that I made with their fathers}

\note{09:32}
\note[item]{The last two weeks were `I will be their God, they shall be my people'}
\note[item]{These phrases are themes throughout the bible}
\note[item]{These phrases are relationship focused}
\note[item]{This week we will discuss the fact that relationships come with a responsibility.}
\note[item]{We're going to strike Lesson 7 because of the gospel meeting}
\note[item]{I'll be out of town for lesson 8 in late April.  Shawn will fill in for me.}
\end{frame}

\begin{goals}
\goal Examine the role that covenants play in God's relationship with man
\goal Establish the importance of faith in both the Old and New Covenants
\goal Reflect on any ways that Christians could `return' to the weaknesses of Old Covenant

\note{09:33}
\note[item]{We'll think about what covenants are.}
\note[item]{We're going to think about all the different covenants God has made with man, and there have been multiple}
\note[item]{Then, we'll think specifically about the Old Covenant also known as the Law of Moses.}
\note[item]{We'll spend a lot of time in Gal. 3 which talks about the purpose of the Old Covenant, the weaknesses of it, and how what we have now is so much better.}
\note[item]{One of the key components of this discussion will be that the Israelites did not pair the adherence to the Law with faith in the Lord.}
\note[item]{We'll talk about how, even though we're not under the Old covenant, how their lack of faith serves as a warning for us.}
\end{goals}

\section{The importance of covenants}

\begin{frame}
\frametitle{A covenant is a contract}
\begin{columns}[T]
\begin{column}{0.5\textwidth}
	\includegraphics[width=\textwidth]{figures/contract.jpg}
\end{column}
\begin{column}{0.5\textwidth}
	Covenants can have\ldots
	\begin{itemize}
	\item requirements
	\item collateral
	\item signatures
	\item payments
	\item dates of effect
	\item etc.
	\end{itemize}
\end{column}
\end{columns}

\note{09:35}
\note[item]{Requirements -- a covenant is really usually a trade}
\note[item]{The government will give my company money if I do work for them.}
\note[item]{Collateral -- there's something given or signed or traded}
\note[item]{Joseph keeps Simeon until the brothers return with Benjamin}
\note[item]{Signatures -- a sign that the contract is ratified}
\note[item]{Boaz traded his shoe to the nearer kinsman to get Ruth.}
\note[item]{Dates of effect -- wills need the death of the testator (Heb 9, lesson 12)}
\note[item]{But, wills rarely make requirements of the people receiving the benefits of the will.  That's not the case with God's covenants.}
\end{frame}

\begin{frame}
\frametitle{Covenants of God through time}
\framesubtitle{}
\begin{columns}[c]
\begin{column}{0.4\textwidth}
	\begin{center}
		`covenant' (ESV)\\
		\includegraphics[width=\columnwidth]{figures/covenant.png}
	\end{center}
	{\footnotesize $-42$ `ark of the covenant'}
\end{column}
\begin{column}{0.6\textwidth}
	\begin{itemize}
		\item{Noah}
		\item{Abraham/Isaac/Jacob}
		\item{Children of Israel}
	\end{itemize}
	\begin{center}
	\includegraphics[width=0.65\columnwidth]{figures/highPriest.jpg}
	\end{center}
	\begin{itemize}
		\item{The New Covenant}
	\end{itemize}

\end{column}
\end{columns}

\note{09:38}
\note[item]{\textbf{Noah}}
\note[item]{Requirements: God saves Noah.  Noah faithful}
\note[item]{Signatures: Rainbow (Gen 8)}
\note[item]{\textbf{Abraham Isaac Jacob}}
\note[item]{Requirements: God promises land, nation, seed.  Abraham walks righteously before God (Gen 17, Lesson 2)}
\note[item]{Signatures: Circumcision}
\note[item]{\textbf{Children of Israel}}
\note[item]{Requirements: God makes Israel succeed physically.  Israelites have to love God but also do a lot of other physical things.}
\note[item]{\textbf{The New Covenant}}
\note[item]{We'll get to that later.}
\end{frame}

\begin{frame}
\frametitle{A primer on Galatians}
\begin{columns}[T]
\begin{column}{0.3\textwidth}
	\includegraphics[width=\columnwidth]{figures/cliffNotes.jpg}
\end{column}
\begin{column}{0.7\textwidth}
	Galatians is a `little' Romans
	\begin{table}
	\begin{tabular}{lll}
~ & Romans & Galatians\\
\hline
The gospel & 1-11 & 1-4 \\ 
Christian living & 12-16 & 5-6 \\
\end{tabular}
\end{table}
\begin{itemize}
	\item Inadequacy of the Law
	\item Justification by Faith
	\item Promises to Abraham
	\item Led by the Spirit
	\item etc.
\end{itemize}
\end{column}
\end{columns}

\note{09:42}
\note[item]{In Romans Paul is trying to remind Christians of their foundation.}
\note[item]{In Galatians is focused on problem of returning to the Old Law (a problem he faced his entire ministry)}
\note[item]{But the both letters end up in the same place.}
\note[item]{First, a long discourse on the what the gospel drawing comparisons to the Abraham's faith and answering questions that followers of the Old Covenant might have}
\note[item]{Second, a discussion of how you apply the gospel to your everyday life}
\note[item]{We're going to use Galatians 3 as our foray into what the Old Covenant was, what it's weaknesses were, and why the children of Israel got it wrong.}
\end{frame}

\begin{frame}
\frametitle{Why was the Old Covenant Given?}
\framesubtitle{Gal. 3:19-26}
\begin{columns}[c]
\begin{column}{0.4\textwidth}
	\includegraphics[width=\columnwidth]{figures/trainingWheels.jpg}
\end{column}
\begin{column}{0.6\textwidth}
	\emph{`added because of transgressions'}\\People tend to fall off the path\\~\\
	\emph{`a mediator is not for just one person, but God is one'}\\ Because Moses gave the Law to the Israelites, they didn't have a first hand relationship with God.\\~\\
	\emph{`was our guardian until Christ'}\\The Old Law constrained people to holy living when they didn't know much about the `plan'.
\end{column}
\end{columns}

\note{09:45}
\note[item]{The Message: The purpose of the law was to keep a sinful people in the way of salvation until Christ (the descendant) came, inheriting the promises and distributing them to us. Obviously this law was not a firsthand encounter with God. It was arranged by angelic messengers through a middleman, Moses. But if there is a middleman as there was at Sinai, then the people are not dealing directly with God, are they? But the original promise is the direct blessing of God, received by faith.}
\note[item]{The Old Covenant was a `guardian' to let people know what holy living looked like, until Jesus came}
\note[item]{But there were problems that came in how people used the Old Law\ldots}
\end{frame}

\section{The importance of faith}

\begin{frame}
\frametitle{All who rely on the works of the law are under a curse}
\framesubtitle{Gal. 3:10-14}
\begin{center}
	\includegraphics[width=\textwidth]{figures/perfectionist.jpg}\\
	Perfect law keeping is not possible.
\end{center}

\note{09:48}
\note[item]{Paul had to constantly combat Judaizing teachers.}
\note[item]{It seems as though they had come in after the fact.}
\end{frame}

\begin{frame}
\frametitle{The `covenant' and the `promise'}
\framesubtitle{Gal. 15-18}
\begin{itemize}
	\item The covenant made with Abraham was binding and came with promises: `land', `nation', and `seed'
	\item The `seed' promise was not fulfilled under the Mosaic Covenant, but under Christ
	\item There was no provision for removal of sins under the Old Covenant, but under the `promise' there was (Hebrews).
	\item There was no provision for Gentiles under the Old Covenant, but under the `promise' there was.
\end{itemize}

\note{09:50}
\note[item]{Paul makes a distinction between covenant made with Abraham and the covenant given to the Israelites.}
\note[item]{Abraham's covenant came with a `land', `nation', and `seed' promise}
\note[item]{The `seed' promise was not fulfilled under the Mosaic Covenant, but under Christ}
\note[item]{This passage, in my mind, makes it clear that the other two promises had already been fulfilled and we were just waiting on Jesus to fulfill the 3rd promise.}
\end{frame}

\begin{frame}
\frametitle{Even the Old Covenant required faith}
\framesubtitle{Gal. 3:6-9}
\begin{columns}[T]
\begin{column}{0.5\textwidth}
	\begin{center}
	\includegraphics[width=\columnwidth]{figures/faith.jpg}\\
	\end{center}
\end{column}
\begin{column}{0.5\textwidth}
	\begin{itemize}
		\item Abraham was justified by faith.
		\item We share Abraham's faith.
		\item And, those saved under the Old Covenant shared Abraham's faith (Heb. 11)
	\end{itemize}
\end{column}
\end{columns}

\note{09:55}
\note[item]{This seems to contradict Galatians 3:12}
\note[item]{Paul is talking about how people tried to \emph{keep} the Law, not the weakness of the Law itself}
\end{frame}

\section{Don't return to the Old Covenant}

\begin{frame}
\frametitle{The Spirit came by faith not the law}
\framesubtitle{Gal. 3:1-5}
\begin{center}
\includegraphics[width=0.8\textwidth]{figures/lemon.jpg}\\
The Galatians were making a bad trade --\\the Old Covenant instead of the New
\end{center}

\note{09:58}
\note[item]{Paul reminds these Christians that they had miraculous gifts among them.}
\note[item]{Obviously, those gifts were greater than anything they had been given under the law of Moses.}
\note[item]{So, he says, why would you want to go back to something that wasn't as good as what you have now?}
\note[item]{Other writers say similar things.  Like the writer of Hebrews.}
\end{frame}

\begin{frame}
\frametitle{`The Law' is both a thing and an idea}
\begin{center}
\includegraphics[width=0.8\textwidth]{figures/law.jpg}\\
The thing: The written Law of Moses\\
The idea: Justification through a system of works
\end{center}

\note{10:00}
\note[item]{We sometimes try to shoehorn every instance of the word `law' into meaning simply the Law of Moses.}
\note[item]{Why do we do that? Because there are many so-called `Christians' who would say that you don't have to do good works to go to heaven.}
\note[item]{These people misunderstand the word `law' to mean that \emph{any} attempt to do good works is futile and has no bearing on whether or not you enter heaven.}
\note[item]{So, to combat that, we say, `Well, these passages are talking about the Law of Moses, not just doing good works in general'.}
\note[item]{But, if we treat the word `law' as meaning simply the Law of Moses, we are taking large chunks of the New Testament and saying, ``That doesn't apply to me.''}
\note[item]{Paul in particularly uses the term `law' to fill a bigger concept.  And, it \emph{does}, I believe, have application for us.}
\end{frame}

\begin{frame}
\frametitle{Don't return to the Old Covenant}
\begin{columns}[c]
\begin{column}{0.35\textwidth}
	\includegraphics[width=\columnwidth]{figures/selfSufficient.jpg}
\end{column}
\begin{column}{0.65\textwidth}
\begin{itemize}
\item The Jews were trying to get closer to God by relying upon their strict adherence to the Law.
\item There was nothing wrong with strict adherence to the Law. (Matt. 23)
\item Clearly, they didn't pair their strict adherence with faith.
\item Are there ways that we can leave faith out of our Christianity?
\end{itemize}
\end{column}
\end{columns}

\note{10:05}
\note[item]{Wasn't faith part of the greatest command and part of the Old Law? Absolutely.}
\note[item]{We've already all the faithful from Hebrews 11, many of lived under the Old Law.}
\note[item]{Where they went wrong was that they wanted to fulfill the outward trappings of the Law (eating only Kosher foods, tithing, etc.), but made no attempt to look beyond their law-keeping to understand the big picture.}
\note[item]{God wanted them to believe Him to have faith in Him, and to look more toward unseen, spiritual things (like heaven) than to physical things}
\note[item]{Isn't that what God wants from us?}
\end{frame}

\section{Review}

\begin{frame}
\frametitle{The covenant I made with their fathers}
	\begin{itemize}
		\item Covenants are the guarantees of relationships
		\item God uses covenants to give promises, and to establish boundaries
		\item The Old Covenant had physical promises, but the spiritual promises came from Abraham
		\item Justification by the Law required perfect law-keeping
		\item Faith was always a requirement for God's people, even under the Old Covenant
		\item We may not try to follow the Old Law, but we still can fall into the trap of trying to be justified by works.
	\end{itemize}

\note{10:12}
\note[item]{Next week we'll talk about God's closeness.  All the ways he showed the Israelites he loved them.}
\note[item]{In two weeks we'll talk about the consequences of breaking God's covenant.}
\end{frame}
